\section{局限性与结论}

\subsection{局限性}
尽管本文方法在纹理细节保真度、
跨视角一致性以及几何对齐方面
取得了较好的效果，
但仍存在一定局限性。
首先，本文方法依赖第一阶段
多视角扩散模型的生成质量。
当多视角结果出现几何错位
或跨视角不一致时，
误差会进一步传递至第二阶段，
从而影响最终纹理质量。
其次，本文仅采用四个正交视角
作为条件输入，
对于顶部、底部以及复杂遮挡区域
的信息覆盖仍然有限。
此外，本文通过二维拼图方式
融合多视角图像特征，
尚未显式建模视角之间的
三维空间关系，
因此在复杂拓扑或高对称结构下，
全局一致性仍有提升空间。
最后，本文方法仍具有较高的
显存与计算开销，
高分辨率纹理生成阶段
对 GPU 资源需求较高，
限制了其实时应用能力。

\subsection{结论}

本文提出了一种多视角引导的
原生三维纹理生成框架MVTrellis，
结合二维多视角扩散模型
与原生三维纹理表示的优势，
实现了高保真、
跨视角一致且几何对齐的
三维纹理生成。

具体而言，
本文首先利用 MV-Adapter
生成与目标几何对齐的
多视角纹理图像，
随后对 Trellis.2 的材质生成阶段
进行多视角条件微调，
使模型能够直接在原生三维空间中
融合多视角外观信息，
从而提升纹理一致性
与不可见区域的补全能力。
实验结果表明，
本文方法在纹理细节、
全局一致性以及复杂区域补全方面
均取得了较好的生成效果，
验证了多视角条件
对于三维纹理生成任务的有效性。

未来工作将进一步探索
更多视角条件、
显式三维一致性建模
以及轻量化实时生成方法，
以提升模型的泛化能力
与实际应用价值。